% =============================================================================
% Preprocessing of the thesis datasets
% File to be \input{} in the main document. Portuguese version: preprocessamento.tex
% Packages used: booktabs, url (or hyperref)
% Keep in sync with scripts/build_thesis_datasets.py.
% Last updated: 2026-11-20 (SIH 2024)
% =============================================================================

\section{Preprocessing}
\label{sec:preprocessing}

Each dataset was preprocessed by a single, reproducible script
(\texttt{scripts/build\_thesis\_datasets.py}), which reads the original files
without modifying them and writes, alongside the final dataset, a step-by-step
exclusion report. Exclusions are sequential: each record is counted only at the
first step in which it is removed.

% -----------------------------------------------------------------------------
\subsection{Hospital Information System (SIH/SUS)}
\label{subsec:pre-sih}

\paragraph{Source and retrieval.}
We used the reduced hospital admission authorization files (RD) of the
Brazilian Unified Health System's Hospital Information System (SIH/SUS),
published by DATASUS at
\url{ftp://ftp.datasus.gov.br/dissemin/publicos/SIHSUS/200801_/Dados}. Each file
covers one federative unit and one \emph{competência}, that is, the month in
which the hospital admission authorization (AIH) was processed for payment. We
downloaded the files for all 27 federative units for the processing months
from January 2024 to June 2025 (486 files) on 19 November 2026, using
\texttt{scripts/download\_datasus\_ftp.py}. For each file, a manifest records the
source path, extraction date, last-modified date on the server, size and
SHA-256 hash, which makes it possible to verify that a rerun uses exactly the
same data.

The files, in DBC format, were decompressed to DBF (\texttt{pyreaddbc}
library), and only 15 of the 113 AIH fields were read: AIH number
(\texttt{N\_AIH}), AIH type (\texttt{IDENT}), admission date (\texttt{DT\_INTER}),
sex (\texttt{SEXO}), race/skin colour (\texttt{RACA\_COR}), age and age unit
(\texttt{IDADE}, \texttt{COD\_IDADE}), principal diagnosis (\texttt{DIAG\_PRINC}),
admission type (\texttt{CAR\_INT}), complexity (\texttt{COMPLEX}), bed specialty
(\texttt{ESPEC}), municipalities of residence and of the hospital
(\texttt{MUNIC\_RES}, \texttt{MUNIC\_MOV}), discharge reason (\texttt{COBRANCA})
and death indicator (\texttt{MORTE}).

\paragraph{Unit of analysis and time frame.}
The unit of analysis is the hospital admission, identified by its AIH. We
selected admissions with an admission date in 2024. Because the processing
month is not the admission month, AIHs for 2024 admissions keep arriving in the
2025 processing months: the 2024 processing months contained 13,335,211 of
these AIHs and the 2025 months another 796,655, of which 510,221 arrived in
January and only 1,118 in June. Reading stopped at the June 2025 processing
month, when new AIHs for 2024 admissions had become negligible (less than
0.01\% of the total). No 2025 admission is part of the dataset: the 2025
processing months were used only to complete the year 2024.

\paragraph{Exclusion criteria.}
Table~\ref{tab:sih-funnel} summarises the exclusions, applied in this order:
\begin{enumerate}
  \item \textbf{Resubmitted AIH.} The same AIH may be processed in more than one
    month (for example, after rejection and resubmission). Only the record from
    the most recent processing month of each \texttt{N\_AIH} was kept.
  \item \textbf{Long-stay continuation AIH} (\texttt{IDENT} = 5). These AIHs
    continue an admission already recorded in a main AIH and would count it
    more than once.
  \item \textbf{Discharge without a defined outcome}, by the first digit of the
    discharge reason (\texttt{COBRANCA}): continued stay (2), in which the
    patient remains in hospital; transfer to another facility (3), in which the
    outcome occurs in another admission; and administrative closure (5).
  \item \textbf{Invalid demographic or diagnosis data}: missing race/skin colour
    (code 99), sex other than male or female, missing age or age above 120
    years, and principal diagnosis outside the ICD-10 chapters. Only 131 2024
    admissions (0.001\%) were excluded at this step, all for age above
    120 years. The records without race/skin
    colour found in the files all belong to long-stay AIHs of admissions that
    started before 2024.
\end{enumerate}

\begin{table}[htbp]
  \centering
  \caption{Exclusions applied to AIHs with an admission date in 2024 (SIH/SUS).}
  \label{tab:sih-funnel}
  \begin{tabular}{lrr}
    \toprule
    Step & Admissions & \% of total \\
    \midrule
    AIHs with admission in 2024 & 14,131,866 & 100.00 \\
    AIH resubmitted in another processing month & $-$96,095 & 0.68 \\
    Long-stay continuation AIH & $-$49,822 & 0.35 \\
    Discharge by continued stay & $-$736,740 & 5.21 \\
    Discharge by transfer & $-$630,287 & 4.46 \\
    Discharge by administrative closure & $-$155,443 & 1.10 \\
    Age above 120 years & $-$131 & 0.00 \\
    Invalid race/skin colour, sex or ICD & 0 & 0.00 \\
    \midrule
    Admissions kept & 12,463,348 & 88.19 \\
    \bottomrule
  \end{tabular}
\end{table}

\paragraph{Outcome.}
The outcome is binary: \emph{death} when the \texttt{MORTE} indicator equals 1,
which covers discharges by death (reason 4) and obstetric deaths (reason 6),
and \emph{discharge} otherwise, which, after the exclusions, corresponds to
regular discharges (reason 1) and obstetric discharges (reason 6). The death
rate in the dataset is 4.98\%. Because it is derived from the discharge record,
the outcome refers to in-hospital death, not to death after discharge.

\paragraph{Protected attributes.}
The protected attributes are race/skin colour (White, Black, Brown, Asian and
Indigenous, following the IBGE categories \emph{Branca}, \emph{Preta},
\emph{Parda}, \emph{Amarela} and \emph{Indígena}), sex (male and female) and age
group (0--17, 18--59 and 60 years or older). Age in years was derived from
\texttt{COD\_IDADE}: ages recorded in days were set to 0 years, ages in months
were converted to completed years, ages in years were kept, and code 5
(hundreds of years) corresponds to $100 + \texttt{IDADE}$. In the SIH, race/skin
colour is recorded by the health facility and is not necessarily
self-reported by the patient, which must be taken into account when
interpreting the fairness metrics.

\paragraph{Other variables.}
Besides the protected attributes and the outcome, the dataset keeps variables
known at admission: age in years, federative unit of the hospital, ICD-10
category of the principal diagnosis (three characters), admission type
(elective, emergency, work accident, traffic accident or other injuries),
complexity (medium or high), bed specialty (surgical, obstetric, clinical,
chronic, psychiatric, sanitary pneumology, paediatric, rehabilitation or
day-hospital/other), admission outside the municipality of residence (when the
six-digit IBGE codes of \texttt{MUNIC\_RES} and \texttt{MUNIC\_MOV} differ) and
admission month. Variables recorded during or after the stay, such as ICU use,
length of stay, procedures performed and amounts paid, were deliberately
excluded to avoid leaking information about the outcome. The principal
diagnosis, although used as an input variable, is finalised when the AIH is
closed, which is a limitation.

\paragraph{Organisation by ICD-10 chapter.}
Because the risk of death and the patient profile vary widely across diagnosis
groups, admissions were split by the 22 ICD-10 chapters, defined by the
official ranges of three-character categories, and each chapter is analysed
separately (Table~\ref{tab:sih-chapters}). Chapters with a death rate below
0.2\% were excluded from the analysis, since their outcome does not vary enough
for the bias analysis, especially in the smaller intersectional groups: XV
(pregnancy, childbirth and the puerperium, 0.04\%), VII (eye and adnexa,
0.02\%) and VIII (ear and mastoid process, 0.20\%\footnote{A rate of 0.197\%,
rounded in the table.}). Chapters with fewer than 1,000 admissions, which do
not allow metrics to be estimated per intersectional group, were also
excluded: XXII (codes for special purposes, 415 admissions). Chapter XX
(external causes, V01--Y98) has no admissions, because its codes are not
accepted as a principal diagnosis in the SIH. Seventeen chapters remain. The
figures for the excluded chapters remain recorded in the exclusion report and
in the dataset summary.

\begin{table}[htbp]
  \centering
  \caption{2024 admissions by ICD-10 chapter of the principal diagnosis (SIH/SUS).}
  \label{tab:sih-chapters}
  \small
  \begin{tabular}{llrrr}
    \toprule
    Ch. & Description (categories) & AIHs in 2024 & Kept & Death (\%) \\
    \midrule
    I     & Infectious and parasitic (A00--B99)                 &   965,970 &   842,452 & 14.95 \\
    II    & Neoplasms (C00--D48)                                & 1,104,508 & 1,002,953 &  8.02 \\
    III   & Blood and immune disorders (D50--D89)               &   138,202 &   120,828 &  5.51 \\
    IV    & Endocrine, nutritional and metabolic (E00--E90)     &   257,186 &   215,324 &  5.95 \\
    V     & Mental and behavioural disorders (F00--F99)         &   339,518 &   183,324 &  0.67 \\
    VI    & Nervous system (G00--G99)                           &   254,320 &   208,990 &  4.64 \\
    VII$^{a}$ & Eye and adnexa (H00--H59)                       &   198,289 &   194,476 &  0.02 \\
    VIII$^{a}$ & Ear and mastoid process (H60--H95)                  &    30,450 &    28,954 &  0.20 \\
    IX    & Circulatory system (I00--I99)                       & 1,332,909 & 1,081,637 &  9.93 \\
    X     & Respiratory system (J00--J99)                       & 1,364,669 & 1,220,759 &  9.91 \\
    XI    & Digestive system (K00--K93)                         & 1,500,519 & 1,397,963 &  3.12 \\
    XII   & Skin and subcutaneous tissue (L00--L99)             &   353,879 &   308,074 &  1.84 \\
    XIII  & Musculoskeletal and connective tissue (M00--M99)    &   294,914 &   271,914 &  0.61 \\
    XIV   & Genitourinary system (N00--N99)                     & 1,073,173 &   990,313 &  3.47 \\
    XV$^{a}$ & Pregnancy, childbirth and puerperium (O00--O99)  & 2,106,315 & 1,999,809 &  0.04 \\
    XVI   & Perinatal conditions (P00--P96)                     &   353,570 &   320,668 &  2.93 \\
    XVII  & Congenital malformations (Q00--Q99)                 &    97,479 &    85,709 &  2.54 \\
    XVIII & Symptoms, signs and abnormal findings (R00--R99)    &   283,376 &   235,614 &  9.79 \\
    XIX   & Injury and poisoning (S00--T98)                     & 1,582,739 & 1,291,275 &  2.56 \\
    XX$^{b}$ & External causes (V01--Y98)                       &         0 &         0 &   --- \\
    XXI   & Contact with health services (Z00--Z99)             &   499,406 &   461,897 &  0.39 \\
    XXII$^{c}$ & Codes for special purposes (U00--U99)          &       475 &       415 &  7.23 \\
    \midrule
          & Total                                               & 14,131,866 & 12,463,348 & 4.98 \\
          & Total of the 17 analysed chapters                   & 11,796,337 & 10,239,694 & 6.05 \\
    \bottomrule
    \multicolumn{5}{l}{\footnotesize $^{a}$ Excluded from the analysis: death rate below 0.2\%.} \\
    \multicolumn{5}{l}{\footnotesize $^{b}$ No admissions: codes not accepted as principal diagnosis in the SIH.} \\
    \multicolumn{5}{l}{\footnotesize $^{c}$ Excluded from the analysis: fewer than 1,000 admissions.} \\
  \end{tabular}
\end{table}

\paragraph{Output files.}
For each analysed chapter, a Parquet file is written
(\texttt{data/processed/thesis/sih\_v2/2024/sih\_v2\_2024\_capNN.parquet}), with
categorical variables stored as categories; together they take about 29~MB.
The dataset comes with the per-chapter exclusion report
(\texttt{sih\_v2\_2024\_exclusoes.csv}), the count of 2024 AIHs per processing
month (\texttt{sih\_v2\_2024\_cobertura\_competencias.csv}) and a per-chapter
summary (\texttt{sih\_v2\_2024\_resumo.parquet}). The dataset can be regenerated
with \texttt{python scripts/download\_datasus\_ftp.py --inicio 2024-01 --fim 2025-06}
followed by \texttt{python scripts/build\_thesis\_datasets.py --only sih}.
